\documentclass[10pt,a4paper]{amsart}
\usepackage[margin=19mm]{geometry}
\usepackage{amsmath,amssymb,mathtools}
\usepackage[scr=rsfs,cal=euler]{mathalfa}
\usepackage{xfrac}
\usepackage[unicode,hidelinks,bookmarks=false]{hyperref}
\newcommand{\ie}{i.e.\ }
\newcommand{\cf}{cf.\ }
\newcommand{\resp}{respectively\ }
\newcommand{\Subsection}{\subsection}
\providecommand{\nobreakdash}{\nobreak\hbox{-}}
\setlength{\emergencystretch}{2em}
\multlinegap0pt
\usepackage{bm}
\usepackage{bbm}
\usepackage{dsfont}
\usepackage{xspace}
\usepackage{cancel}
\usepackage{mathtools}
\usepackage{amsthm}
\makeatletter
\@ifundefined{theorem}{\newtheorem{theorem}{Theorem}}{}
\@ifundefined{proposition}{\newtheorem{proposition}[theorem]{Proposition}}{}
\makeatother
\title[Novel periodic solutions and rogue waves of the defocusing s]{Novel periodic solutions and rogue waves of the defocusing scalar and coupled Ablowitz-Ladik systems on a nonzero background}
\author{Francesco Coppini, Barbara Prinari}
\date{Source: arXiv:2606.05123v1 (CC BY 4.0)}
\begin{document}
\maketitle

\noindent\emph{Source and licence.} Novel periodic solutions and rogue waves of the defocusing scalar and coupled Ablowitz-Ladik systems on a nonzero background. Version arXiv:2606.05123v1 (CC BY 4.0), distributed under the Creative Commons Attribution 4.0 International licence, as recorded in the arXiv licence metadata for this exact version and shown on the abstract page https://arxiv.org/abs/2606.05123v1. Full rights notice: licenses/arxiv-2606.05123-CC-BY-4.0.txt.\par\smallskip

\noindent\textit{Editorial scope.} Counted unit: the Proposition whose source label is \texttt{prop:main} in the article (source0.tex lines 1379--1406, source label \texttt{prop:main}), delivered together with the article's own complete original proof of it (source0.tex lines 1408--1482, one \texttt{proof} environment of 75 lines). No mathematics is written, repaired or re-derived: statement and proof are the article's own text line for line. Required context retained uncounted: none. Every definition, display and label of those lines is retained unchanged. Omitted from this package and not redistributed: 1--1378, 1407--1407, 1483--1507. Nothing inside a retained excerpt is omitted or altered: every retained body is delivered exactly as the article's lines are written, so no omission receipt of any kind accompanies this package. Citations: the retained excerpts carry no citation command at all, so no bibliography entry of the article is transcribed and no reference is redistributed. Reference adaptation: none was needed and none was made; every reference inside the delivered unit resolves inside it, so no unresolved reference or citation is produced. Printed slips kept literal: no printed word, symbol, hypothesis or number of the counted statement or of its proof was altered. Layout and class: the article's own class is \texttt{article} with packages including geometry, graphicx, epsfig, amsmath, amssymb, amsfonts, xcolor, cite; the wrapper is the lane's pinned \texttt{amsart} editorial wrapper at 10pt on A4, which restates the theorem-like environments the retained text needs and adds the article's own macro definitions that the retained text needs (none), with extra wrapper packages (bm, bbm, dsfont, xspace, cancel, mathtools). The delivered numbering is the unit's own. Assets: no retained span contains a figure environment, a graphics-inclusion command, a PDF-page inclusion, a table or a TikZ picture; no image file is copied into this package, the package declares no asset dependency and no image file is redistributed with it. Licence: version arXiv:2606.05123v1 of this article is distributed under the Creative Commons Attribution 4.0 International licence, as recorded in the arXiv licence metadata for this exact version and shown on the abstract page https://arxiv.org/abs/2606.05123v1. A full rights notice is delivered at licenses/arxiv-2606.05123-CC-BY-4.0.txt. Characters: every retained excerpt is pure ASCII, so the delivered engine, XeLaTeX with its default Latin Modern text font, reports no missing character.
\begin{proposition}\label{prop:main}
Let $\xi_0=k/2$.  With $s:=\sinh(k/2)$, $A:=\cosh(k/2)$, and
$Q:=\sqrt{s^2-\rho^2 A^2}$ (positive by the KM condition), the
maximum modulus of the regular KM breather is given \emph{exactly} by
\begin{equation}\label{eq:maxmod}
  \max_{n\in\mathbb{Z},\;t\in\mathbb{R}} |q_n(t)|
  \;=\;
  \rho\,\frac{C_1+s}{s-C_2}
  \;=\;
  \rho\,\frac{A\bigl(s\sqrt{1-\rho^2}+Q\bigr)-2\rho A}{sQ-\rho A^2/A}
  \;=\;
  \frac{\rho\bigl[As\sqrt{1-\rho^2}-2\rho A+sQ\bigr]}{sQ-\rho A},
\end{equation}
where $C_1:=A(\cosh\phi-2)/\sinh\phi$.  Furthermore,
\begin{equation}\label{eq:crit}
  \max_{n\in\mathbb{Z},\;t\in\mathbb{R}}|q_n(t)| < 1
  \quad\Longleftrightarrow\quad
  \rho < \tfrac{1}{2}
  \;\text{ and }\;
  k > k_c(\rho),
\end{equation}
where $k_c(\rho)$ is the unique solution $k>2\operatorname{arctanh}\rho$
of the implicit equation
\begin{equation}\label{eq:kstar}
  A\bigl(s\sqrt{1-\rho^2}+1-2\rho\bigr) = s(1-\rho)\,\frac{Q}{\rho},
  \qquad s=\sinh\tfrac{k}{2},\; A=\cosh\tfrac{k}{2},\; Q=\sqrt{s^2-\rho^2 A^2}.
\end{equation}
\end{proposition}

\begin{proof}
\textbf{Step 1: reduction to $n=0$.}
Writing $|q_n|^2 = \rho^2(N_{\mathrm{re}}^2+N_{\mathrm{im}}^2)/D_n^2$ with
$N_{\mathrm{re}}=S_n+C_1\cos\theta$, $N_{\mathrm{im}}=A\sin\theta$, one sees
that the ratio $|q_n|^2$ is maximized at the site $n=0$ for all valid
$(\rho,k)$.  With $\xi_0=k/2$ we have $S_0=-s$ and $S_n=\pm\sinh(k|n-\tfrac{1}{2}|)$
grows exponentially with $|n|$, making $|q_n|\to\rho$ exponentially fast
away from $n=0$.
 
\textbf{Step 2: exact maximum at $n=0$.}
Setting $u:=\cos\theta\in[-1,1]$, the modulus squared at $n=0$ is
\[
  |q_0|^2 = \rho^2\,f(u), \qquad
  f(u) = \frac{(C_1^2-A^2)u^2 - 2sC_1 u + s^2+A^2}{(s-C_2 u)^2},
\]
where $S_0=-s$ and $D_0=-(s-C_2 u)>0$ (denominator sign fixed by regularity).
Setting $f'(u)=0$ gives the unique interior critical point
\[
  u^* = \frac{s^2(C_1+C_2)+C_2 A^2}{s(C_1^2-A^2+C_1 C_2)}.
\]
One computes $C_1^2-A^2+C_1C_2 = -3A^2(\cosh\phi-1)/\sinh^2\phi < 0$
and $s^2(C_1+C_2)+C_2 A^2 > 0$, so $u^*<0$.  A further estimate shows
$u^*<-1$, i.e.\ $u^*\notin[-1,1]$, so $f$ is monotone on $[-1,1]$ and
its maximum is at $u=-1$.
 
\textbf{Step 3: closed form.}
Evaluating at $u=-1$:
\[
  f(-1) = \frac{(C_1^2-A^2)+2sC_1+s^2+A^2}{(s-C_2)^2}
        = \frac{(C_1+s)^2}{(s-C_2)^2},
\]
so
\[
  \max_{t}|q_0(t)| = \rho\,\frac{C_1+s}{s-C_2}
\]
(both numerator and denominator are positive: $C_1+s>0$ since $C_1>-s$
for all KM parameters, and $s-C_2>0$ by the regularity condition
$s\sinh\phi>A$, i.e.\ $s>C_2$).
Substituting $C_1=A(\delta-2)/\sinh\phi$, $C_2=A/\sinh\phi$,
$\delta=\sqrt{1-\rho^2}\,s/\rho$, and $\sinh\phi=Q/\rho$
yields the formula in~\eqref{eq:maxmod}.
 
\textbf{Step 4: condition $\max|q_n|<1$.}
The inequality $\rho(C_1+s)<s-C_2$ (equivalently $\rho f(-1)<1$) expands,
after substitution of $C_1$, $C_2$ in terms of $s,A,Q,\rho$, to
\[
  \rho\bigl[As\sqrt{1-\rho^2}-2\rho A+sQ\bigr] < sQ-\rho A.
\]
Rearranging:
\[
  \rho A\bigl(s\sqrt{1-\rho^2}+1-2\rho\bigr) < sQ(1-\rho).
\]
For $\rho<1/2$ the factor $(s\sqrt{1-\rho^2}+1-2\rho)$ is always positive
(since $s>0$ and $1-2\rho>0$), and both sides are positive, so the
inequality is strict and is equivalent to~\eqref{eq:kstar} being reversed.
As $k\to\infty$ the left side grows as $\rho As\sqrt{1-\rho^2}\sim\rho s^2\sqrt{1-\rho^2}$
while $sQ\sim s^2$, so the right side dominates and the condition holds for
all sufficiently large $k$.
Conversely, for $k$ just above the KM threshold $Q\to 0$ and the right side
vanishes, so the condition fails.  By continuity there is a unique crossing
$k=k_c(\rho)$ above which $\max|q_n|<1$.
 
For $\rho=1/2$: the condition becomes $A(s/\sqrt{3}+0)<sQ$, i.e.\ $A/\sqrt{3}<Q$,
while the KM condition requires $Q>0$, i.e.\ $s>A/2=s_{\mathrm{KM}}$.
At $k=k_c$ one needs simultaneously $k>k_{\mathrm{KM}}$ and the
equality in~\eqref{eq:kstar}; these two conditions reduce to $s=A/\sqrt{3}$,
which is incompatible with $s>\rho A=A/2$ (required by KM) since
$1/\sqrt{3}<1/2$ is false.  Hence no valid $k$ exists and $\max|q_n|\geq1$
for all valid parameters.
 
For $\rho>1/2$: the factor $1-2\rho<0$, so the left side of the
rearranged inequality can be negative for small $s$, but checking the
large-$k$ limit gives $\max|q_n|\to 2\rho>1$, so the condition fails
for all $k$.
\end{proof}

\end{document}
